\documentclass[11pt]{article}
\usepackage[utf8]{inputenc}
\usepackage[T1]{fontenc}
\usepackage{geometry}
\usepackage{booktabs}
\usepackage{longtable}
\usepackage{array}
\usepackage{float}
\usepackage{hyperref}
\usepackage{enumitem}
\usepackage{graphicx}
\geometry{margin=1in}
\hypersetup{colorlinks=true, linkcolor=blue, urlcolor=blue}

\title{Sound of Culture\\Phase 1.2 Report: Ultimate Guitar Pipeline and Final Dataset Construction}
\author{Detailed project report}
\date{\today}

\begin{document}
\maketitle
\tableofcontents
\clearpage

\section{Purpose of this report}
This report documents the retained song-side construction workflow that starts
from the operational country-only artist file, runs Ultimate Guitar discovery
and download, supplements the missing chart material, and ends with the final
country-only chord dataset and the recovery of the remaining missing metadata.

The key distinction relative to Phase 1.1 is that this report is about
\emph{song-side} construction rather than \emph{artist-side} construction. The
country-only file inherited from Phase 1.1 is treated here as the operational
artist input. The report therefore focuses on:
\begin{enumerate}[leftmargin=2em]
\item Ultimate Guitar discovery, planning, and raw JSON download;
\item Billboard supplementation for still-missing artists or songs;
\item final dataset assembly on top of \texttt{Restricted Final v6};
\item the hierarchical recovery of \texttt{release\_year} and other residual
metadata;
\item snapshot-assisted replication of the retained final outputs.
\end{enumerate}

\section{Inputs inherited from Phase 1.1}
The song-side pipeline does not start from the full artist universe. It starts
from the operational \texttt{artist\_universe\_country\_only.csv}. This is a
retained design choice. The project does \emph{not} use the adjacent-only
universe as the main scrape frame for the final country-only dataset.

That distinction is essential:
\begin{itemize}[leftmargin=2em]
\item \texttt{country-only} is the core artist frame used to search Ultimate
Guitar and to attach artist metadata to the final dataset;
\item \texttt{adjacent-only} is retained as a broader documented perimeter, not
as the main song-discovery denominator for this report's dataset.
\end{itemize}

\section{Retained pipeline}

\subsection{Stage 1. Ultimate Guitar discovery, planning, and raw JSON download}
The retained Ultimate Guitar pipeline is implemented in
\texttt{code/country\_only\_ultimate\_guitar.py}. The main retained entrypoints
are:
\begin{itemize}[leftmargin=2em]
\item \texttt{run\_discovery()} at lines 297--348;
\item \texttt{build\_plan()} at lines 375--417;
\item \texttt{run\_download()} and its worker logic at lines 420--560.
\end{itemize}

The script is intentionally split into three commands:
\begin{enumerate}[leftmargin=2em]
\item \texttt{discover}, which searches the country-only artist list and stores
up to the top three versions per song;
\item \texttt{plan}, which compares discovery targets with the raw JSON files
already on disk;
\item \texttt{download}, which retrieves the remaining JSON payloads with
resume logic, sharding options, and explicit handling of rate limits and legal
451 responses.
\end{enumerate}

This separation is retained because it makes the workflow restartable and
auditable. Discovery, planning, and downloading are conceptually distinct
objects and should remain distinct operationally.

\begin{table}[H]
\centering
\caption{Retained Ultimate Guitar plan snapshot}
\begin{tabular}{p{0.34\textwidth}rrr}
\toprule
Layer & Songs & JSON total & JSON missing \\
\midrule
UG chords & 55,335 & 64,829 & 11 \\
UG tabs & 5,243 & 7,299 & 3 \\
Combined & --- & 72,128 & 14 \\
\bottomrule
\end{tabular}
\end{table}

The important substantive implication is that the pipeline keeps multiple JSON
versions per song when available. The retained workflow is not a crude
top-result-only scrape.

\subsection{Stage 2. Billboard supplement}
The retained supplement logic is implemented in
\texttt{code/supplement\_billboard\_country\_chords.py}, especially:
\begin{itemize}[leftmargin=2em]
\item \texttt{run()} at lines 235--363;
\item \texttt{main()} at lines 366--390.
\end{itemize}

This stage exists because broad Ultimate Guitar discovery still leaves two
types of holes:
\begin{itemize}[leftmargin=2em]
\item artists missing from the original country-only scrape frame but present
in the Billboard repair layer;
\item songs by already-known artists that remain absent from the discovery
output.
\end{itemize}

The retained supplement therefore has two branches:
\begin{enumerate}[leftmargin=2em]
\item broad discovery for newly added artists when they are still missing from
the discovery output;
\item targeted song search for the Billboard songs that remain missing after
the broad pass.
\end{enumerate}

\begin{table}[H]
\centering
\caption{Saved Billboard supplement report used by the retained workflow}
\begin{tabular}{p{0.52\textwidth}r}
\toprule
Metric & Value \\
\midrule
Added artists in supplement & 206 \\
Target songs processed in the saved run report & 1,516 \\
Existing discovery rows before supplement & 70,496 \\
Rows added by the supplement & 13 \\
Final discovery rows after merge & 70,504 \\
New raw JSON downloads & 8 \\
Attempted target keys recorded & 1,087 \\
\bottomrule
\end{tabular}
\end{table}

This is also the stage that explains why the current operational
\texttt{artist\_universe\_country\_only.csv} is larger than the clean Phase 1.1
core. The artist frame used in the final song pipeline is intentionally
expanded when Billboard recovery requires it.

\subsection{Stage 3. Build the final chord-level dataset}
The retained final builder is
\texttt{code/build\_country\_only\_chords\_final.py}. The main retained blocks
are:
\begin{itemize}[leftmargin=2em]
\item \texttt{build\_new\_song\_rows()} at lines 595--643;
\item \texttt{apply\_num\_json\_used()} at lines 686--705;
\item \texttt{apply\_chart\_variables()} at lines 1919--1953;
\item \texttt{build\_final\_dataset()} at lines 1956--1983;
\item \texttt{main()} at lines 2015--2033.
\end{itemize}

The builder starts from \texttt{Sound\_of\_Culture\_Country\_Restricted\_Final\_v6}
and extends it. The retained logic is:
\begin{enumerate}[leftmargin=2em]
\item load the discovery rows and the raw JSON payloads;
\item turn newly available songs into song-level rows;
\item deduplicate those new rows first by Ultimate Guitar ID and then by
artist-song key;
\item reattach artist metadata from the operational country-only file;
\item create bridge rows for artists who have raw JSON but no existing row in
the combined file;
\item record how many chord JSON files were actually used for each song;
\item attach Billboard year-end variables;
\item recover missing \texttt{release\_year};
\item save the final \texttt{csv} and \texttt{dta}.
\end{enumerate}

\begin{table}[H]
\centering
\caption{Retained final output counts}
\begin{tabular}{p{0.56\textwidth}r}
\toprule
Object & Count \\
\midrule
Current new-song cache (\texttt{country\_only\_chords\_new\_songs\_2026\_04\_08.csv}) & 36,584 \\
Current final country-only chord dataset & 44,058 \\
\bottomrule
\end{tabular}
\end{table}

Two methodological details are retained and should not be misunderstood.

First, \texttt{genre} in the final dataset is a \emph{song-level} attribute
read from the Ultimate Guitar payload. It is not imposed from the artist-side
country-only classification. This is why songs by country-only artists can
still appear with non-country song tags.

Second, the builder can use up to three JSON payloads for the same song. The
retained workflow deliberately prefers deeper recovery from the best available
versions rather than hard-coding a top-1 rule.

\subsection{Stage 4. Recover missing \texttt{release\_year} and residual metadata}
The retained \texttt{release\_year} logic is in
\texttt{code/build\_country\_only\_chords\_final.py}, especially
\texttt{enrich\_release\_years()} at lines 1519--1679.

The retained hierarchy is explicit:
\begin{enumerate}[leftmargin=2em]
\item manual web-confirmed overrides;
\item historical master caches inherited from previous project layers;
\item dedicated caches saved by the current country-only final run;
\item artist-level iTunes lookup;
\item artist-level MusicBrainz recording lookup;
\item artist-level MusicBrainz release lookup;
\item artist-level Deezer lookup;
\item song-bundle iTunes lookup for residual missing songs;
\item song-level fallback through Discogs, MusicBrainz, Wikidata, and finally
Wikipedia.
\end{enumerate}

This is a retained design choice because it minimizes both cost and noise. The
workflow first reuses already-confirmed knowledge, then queries by artist when
many missing songs belong to the same performer, and only at the end falls back
to song-by-song searches.

\begin{table}[H]
\centering
\caption{Saved augmentation and residual-missingness diagnostics used by the final pipeline}
\begin{tabular}{p{0.56\textwidth}r}
\toprule
Metric & Value \\
\midrule
Base rows in the saved Wikipedia year-end supplement summary & 2,729 \\
Supplemental rows added from the year-end supplement summary & 473 \\
Augmented missing-song targets in the retained v4 summary & 1,515 \\
Targets linked to existing artists in the retained v4 summary & 1,490 \\
Targets still missing an artist in the retained v4 summary & 25 \\
Rows in the retained residual release-year file & 2,893 \\
\bottomrule
\end{tabular}
\end{table}

The key interpretation is that missing-data recovery is not a separate afterthought.
It is part of the final build itself and is handled by a staged hierarchy rather
than a single source.

\subsection{Stage 5. Snapshot-assisted replication and package refresh}
The retained workflow also preserves the package layer that restores or refreshes
the canonical outputs. The main scripts are:
\begin{itemize}[leftmargin=2em]
\item \texttt{code/run\_country\_songs\_replication.py}, especially lines
10--29, 67--77, 87--108, 111--170;
\item \texttt{code/run\_country\_merge\_v6\_replication.py}, especially lines
67--71, 109--176, 215--247, 294--360.
\end{itemize}

These wrappers are retained because they let the project:
\begin{itemize}[leftmargin=2em]
\item restore the canonical upstream song and merge-to-v6 layers quickly in a
clean workspace;
\item refresh the packaged snapshots after a validated retained run;
\item distinguish between exact reproducibility of final deliverables and a
full live rebuild from raw web sources.
\end{itemize}

\section{Retained versus non-retained choices}
\subsection{Retained choices}
The following choices are retained in the final workflow:
\begin{itemize}[leftmargin=2em]
\item using the country-only file, not the adjacent universe, as the main
Ultimate Guitar scrape input;
\item separating discovery, planning, and downloading into distinct commands;
\item using a targeted Billboard supplement instead of inflating the base
discovery indiscriminately;
\item constructing the final dataset from \texttt{Restricted Final v6} plus new
UG rows, rather than replacing the upstream file wholesale;
\item using a hierarchical cache-first \texttt{release\_year} recovery design;
\item preserving exact-replication wrappers for the final deliverables.
\end{itemize}

\subsection{Not retained as final choices}
The following ideas are not retained as the final workflow:
\begin{itemize}[leftmargin=2em]
\item scraping the full adjacent artist universe before the country-only song
pipeline is stabilized;
\item forcing artist-level genre identity onto song-level \texttt{genre} values
in the final dataset;
\item treating snapshot-assisted replication as if it were equivalent to the
full live historical build.
\end{itemize}

Those choices were rejected because they would either make the final country
dataset less interpretable or blur the difference between provenance and exact
restoration.

\section{Exact scripts used in the retained workflow}
\begin{longtable}{p{0.36\textwidth}p{0.18\textwidth}p{0.33\textwidth}}
\toprule
Script copy in this folder & Main line ranges used & Role in the retained workflow \\
\midrule
\endfirsthead
\toprule
Script copy in this folder & Main line ranges used & Role in the retained workflow \\
\midrule
\endhead
\texttt{code/country\_only\_ultimate\_guitar.py} & 297--348; 375--417; 420--560 & Discovery, planning, and raw JSON download for the operational country-only artist frame. \\
\texttt{code/supplement\_billboard\_country\_chords.py} & 235--363; 366--390 & Billboard-based recovery of added artists and still-missing songs. \\
\texttt{code/build\_country\_only\_chords\_final.py} & 595--643; 686--705; 1519--1679; 1919--1983; 2015--2033 & Build new song rows, attach chart variables, recover release years, and save the final dataset. \\
\texttt{code/run\_country\_songs\_replication.py} & 10--29; 67--77; 87--108; 111--170 & Restore and refresh the upstream country-song package outputs and caches. \\
\texttt{code/run\_country\_merge\_v6\_replication.py} & 67--71; 109--176; 215--247; 294--360 & Restore and verify the merge-to-v6 bridge that the final builder depends on. \\
\bottomrule
\end{longtable}

\section{Stata note}
There is no retained substantive Stata construction workflow in Phase 1.2. The
main pipeline is Python-based. Stata appears only as a thin launcher layer for
replication packages, so this report does not include Stata code snippets.

\section{Conclusion}
Phase 1.2 turns the artist-side frame into the final song-side construction
dataset. The retained logic is a layered one:
\begin{itemize}[leftmargin=2em]
\item discover and download Ultimate Guitar material from the country-only
artist frame;
\item repair the coverage gaps that remain through the Billboard supplement;
\item append the valid new songs to the upstream v6 base;
\item recover \texttt{release\_year} through a documented hierarchy;
\item preserve exact restoration wrappers for the canonical final outputs.
\end{itemize}

That retained sequence is the final Phase 1.2 construction workflow used by the
project.

\end{document}
